\uuid{avc0}
\exo7id{1164}
\titre{exo7 1164}
\auteur{barraud}
\organisation{exo7}
\datecreate{2003-09-01}
\isIndication{false}
\isCorrection{false}
\chapitre{Déterminant, système linéaire}
\sousChapitre{Système linéaire, rang}
\cpge{mpsi-mp2i.calcul-matriciel-systemes.c, pcsi.calcul-matriciel-systemes-lineaires.c, ptsi.calcul-matriciel-systemes-lineaires.c, mpsi-mp2i.matrices-applications-lineaires.c, pcsi.matrices-applications-lineaires.d, ptsi.matrices-applications-lineaires.d, bcpst1.systemes-lineaires.a, bcpst1.systemes-lineaires.c}
\module{Algèbre}
\niveau{L2}
\difficulte{}

\contenu{
\question{
Sans chercher à résoudre les systêmes suivants, discuter la nature
  de leurs ensembles de solution :
  $$
  \left\{
    \begin{array}{rcrcrcl}
         x &+&    y &-&    z &=&  0  \\
         x &-&    y & &      &=&  0  \\
         x &+&    y &+&    z &=&  0  
    \end{array}
  \right.  \qquad \left\{
    \begin{array}{rcrcrcl}
         x &+&  3 y &+&  2 z &=&  1  \\
       2 x &-&  2 y & &      &=&  2  \\
         x &+&    y &+&    z &=&  2  
    \end{array}
  \right.  \qquad \left\{
    \begin{array}{rcrcrcl}
         x &+&  3 y &+&  2 z &=&  1  \\
       2 x &-&  2 y & &      &=&  2  \\
         x &+&    y &+&    z &=&  3   
    \end{array}
  \right.
  $$
}
}
